% ============================================================
% CHƯƠNG 2: PHÂN TÍCH THIẾT KẾ HỆ THỐNG
% ============================================================
\chapter{PHÂN TÍCH THIẾT KẾ HỆ THỐNG}

Chương này trình bày toàn bộ quá trình phân tích và thiết kế hệ thống. Nội dung bao gồm: phân tích yêu cầu chức năng và phi chức năng, thiết kế kiến trúc tổng thể theo mô hình dịch vụ vi mô với sáu dịch vụ độc lập, thiết kế chi tiết Dịch vụ Tác tử -- trái tim AI của hệ thống với đồ thị trạng thái LangGraph 14 nút và sáu tác tử chuyên biệt, thiết kế cơ sở dữ liệu quan hệ kết hợp véc-tơ, chiến lược tìm kiếm lai và hệ thống giám sát toàn diện.

\section{Phân tích yêu cầu}

\subsection{Yêu cầu chức năng}

Hệ thống phục vụ ba nhóm đối tượng chính: người dùng cuối (người mua, người thuê, nhà đầu tư), quản trị viên hệ thống, và bản thân hệ thống (các tác vụ tự động). Dưới đây là các nhóm chức năng được phân tích chi tiết.

\subsubsection{Nhóm chức năng dành cho người dùng cuối}

Đây là nhóm chức năng cốt lõi, quyết định trải nghiệm của người dùng khi tương tác với nền tảng:

\begin{enumerate}
    \item \textbf{Xác thực và quản lý tài khoản:}
    \begin{itemize}
        \item Đăng ký tài khoản mới với email và mật khẩu. Mật khẩu được hash bằng bcrypt trước khi lưu trữ.
        \item Đăng nhập bằng email/mật khẩu, hệ thống trả về JWT token có thời hạn 24 giờ.
        \item Người dùng chưa đăng nhập (anonymous) vẫn có thể sử dụng chatbot với giới hạn 20 lượt/ngày. Người dùng đã xác thực được nâng lên 200 lượt/ngày.
        \item Hồ sơ người dùng bao gồm: họ tên, email, số điện thoại, ảnh đại diện, quyền (is\_admin, is\_superuser).
    \end{itemize}

    \item \textbf{Duyệt và tìm kiếm bất động sản:}
    \begin{itemize}
        \item \textbf{Danh sách bất động sản bán:} Hiển thị dưới dạng lưới thẻ (card grid) với phân trang. Mỗi thẻ hiển thị: tiêu đề, giá, diện tích, địa chỉ (quận, thành phố), số phòng ngủ, loại tin.
        \item \textbf{Danh sách bất động sản cho thuê:} Tương tự như trang bán nhưng dành riêng cho tin cho thuê.
        \item \textbf{Bộ lọc nâng cao:} Người dùng có thể lọc theo 10+ tiêu chí: loại tin (bán/cho thuê), loại hình bất động sản (căn hộ, nhà riêng, đất nền, biệt thự,\dots), thành phố, quận/huyện, khoảng giá, diện tích, số phòng ngủ, số phòng tắm, hướng nhà.
        \item \textbf{Tìm kiếm theo từ khóa:} Hỗ trợ tìm kiếm full-text trên các trường tiêu đề, mô tả, địa chỉ, quận, thành phố.
        \item \textbf{Sắp xếp kết quả:} Theo giá tăng dần/giảm dần, diện tích tăng dần/giảm dần, mới nhất.
    \end{itemize}

    \item \textbf{Xem chi tiết bất động sản:}
    \begin{itemize}
        \item Trang chi tiết hiển thị đầy đủ thông tin của một bất động sản: tiêu đề, mô tả chi tiết, giá (dạng số + text), đơn giá theo m\textsuperscript{2}, diện tích, địa chỉ đầy đủ (phường/xã, quận/huyện, tỉnh/thành phố), tọa độ địa lý (latitude, longitude).
        \item Thông số kỹ thuật: số phòng ngủ, số phòng tắm, số tầng, hướng nhà, mặt tiền, đường vào.
        \item Thông tin pháp lý và nội thất: tình trạng pháp lý (sổ đỏ, sổ hồng,\dots), tình trạng nội thất (đầy đủ, cơ bản, bàn giao thô).
        \item \textbf{Bất động sản tương tự:} Hệ thống tự động gợi ý các bất động sản có đặc điểm tương đồng (cùng khu vực, cùng phân khúc giá, cùng loại hình).
    \end{itemize}

    \item \textbf{Thống kê và phân tích thị trường:}
    \begin{itemize}
        \item \textbf{Tổng quan thị trường:} Hiển thị các chỉ số chính: tổng số tin đăng đang hoạt động, giá trung bình, diện tích trung bình, số lượng tin bán và cho thuê.
        \item \textbf{Top khu vực:} Bảng xếp hạng các thành phố/quận có nhiều tin đăng nhất, có thể lọc theo loại tin.
        \item \textbf{Thống kê giá theo quận:} Với mỗi quận trong một thành phố, hiển thị số lượng tin, giá trung bình, giá thấp nhất, giá cao nhất và đơn giá trung bình/m\textsuperscript{2}.
        \item \textbf{Phân bố loại hình:} Thống kê số lượng tin theo từng loại hình bất động sản.
    \end{itemize}

    \item \textbf{Tương tác với chatbot tư vấn đa tác tử:}
    \begin{itemize}
        \item \textbf{Chat widget nổi:} Luôn hiển thị ở góc phải dưới màn hình, có thể mở rộng/thu gọn.
        \item \textbf{Hội thoại đa lượt:} Hệ thống lưu trữ toàn bộ lịch sử hội thoại theo phiên (session). Người dùng có thể xem lại các phiên chat trước đó.
        \item \textbf{Phản hồi có trích dẫn nguồn:} Mỗi câu trả lời của chatbot đều kèm danh sách nguồn tham khảo (tên bất động sản, đường dẫn, điểm liên quan), cho phép người dùng kiểm chứng thông tin.
        \item \textbf{Gợi ý câu hỏi tiếp theo:} Sau mỗi phản hồi, chatbot đề xuất 2--3 câu hỏi liên quan.
        \item \textbf{Đánh giá chất lượng:} Người dùng có thể đánh giá phản hồi (thích/không thích) và để lại nhận xét, giúp cải thiện chất lượng hệ thống.
        \item \textbf{Cá nhân hóa:} Hệ thống ghi nhớ sở thích của người dùng (khu vực ưa thích, ngân sách, loại hình) và sử dụng chúng để cá nhân hóa phản hồi trong các phiên sau.
    \end{itemize}
\end{enumerate}

\subsubsection{Nhóm chức năng chatbot đa tác tử -- Chi tiết}

Đây là nhóm chức năng cốt lõi tạo nên sự khác biệt của hệ thống so với các nền tảng bất động sản truyền thống. Hệ thống triển khai kiến trúc hai lớp: đường ống nội bộ trong Backend xử lý nhanh các truy vấn đơn giản, và Dịch vụ Tác tử (Agent Service) chạy đồ thị trạng thái LangGraph riêng biệt cho suy luận đa bước nâng cao.

\begin{enumerate}
    \item \textbf{Phân loại ý định người dùng (Intent Routing):}
    \begin{itemize}
        \item Hệ thống tự động phân tích câu hỏi tiếng Việt để xác định người dùng đang cần: tìm kiếm bất động sản, phân tích thị trường, tư vấn pháp lý, tư vấn đầu tư, tra cứu dự án, hay cập nhật tin tức.
        \item \textbf{Ba cơ chế routing:} (1) Keyword-based -- sử dụng regex và từ khóa tiếng Việt không dấu, hoạt động không cần API key; (2) LLM-based -- sử dụng Google Gemini 2.0 Flash để phân loại với độ chính xác cao hơn; (3) Hybrid -- kết hợp keyword cho tốc độ và LLM cho độ chính xác.
        \item Router có thể chọn nhiều agent cùng lúc nếu câu hỏi liên quan đến nhiều lĩnh vực.
    \end{itemize}

    \item \textbf{Sáu tác tử chuyên biệt hoạt động song song:}
    \begin{itemize}
        \item \textbf{Property Search Agent:} Tìm kiếm bất động sản theo yêu cầu -- tự động trích xuất bộ lọc từ truy vấn (giá, diện tích, khu vực, loại hình, số phòng), thực hiện hybrid search kết hợp SQL filtering và pgvector semantic search, xếp hạng lại qua Cohere Rerank.
        \item \textbf{Investment Advisor Agent:} Phân tích đầu tư -- đánh giá ROI, dòng tiền cho thuê, tỷ suất vốn hóa dựa trên dữ liệu thực tế từ thị trường, kèm theo khuyến cáo tài chính bắt buộc.
        \item \textbf{Market Analysis Agent:} Phân tích thị trường -- cung cấp số liệu thống kê giá theo quận, xu hướng thị trường, snapshot giá trung bình (không phải chuỗi thời gian).
        \item \textbf{News Agent:} Tra cứu tin tức -- tìm kiếm trong kho bài viết tin tức thị trường mới nhất.
        \item \textbf{Project Agent:} Tra cứu dự án -- tìm kiếm thông tin chủ đầu tư, quy mô, vị trí, tiện ích của các dự án bất động sản.
        \item \textbf{Legal Advisor Agent:} Tư vấn pháp lý -- truy xuất từ kho kiến thức pháp lý (sổ đỏ, thủ tục sang tên, thuế phí, quy hoạch), kèm theo khuyến cáo pháp lý bắt buộc.
    \end{itemize}

    \item \textbf{Suy luận đầu tư chuyên sâu (Investment Model):}
    \begin{itemize}
        \item Khi người dùng hỏi về đầu tư, hệ thống xây dựng mô hình tài chính với các chỉ số: giá bất động sản, dòng tiền cho thuê kỳ vọng, tỷ lệ vốn vay, lãi suất, kỳ hạn vay, chi phí vận hành, tỷ lệ trống.
        \item Tính toán các chỉ số đầu tư: ROI, NPV (Net Present Value), dòng tiền hàng năm, điểm hòa vốn.
    \end{itemize}

    \item \textbf{Đánh giá hội đồng (Committee Review):}
    \begin{itemize}
        \item Kết quả phân tích đầu tư được xem xét từ ba góc nhìn độc lập: bullish (lạc quan), bearish (thận trọng), và skeptical (hoài nghi).
        \item Mỗi góc nhìn đưa ra lập luận, mức độ rủi ro và hành động đề xuất dựa trên cùng bộ bằng chứng.
    \end{itemize}

    \item \textbf{Cơ chế ReAct (Reasoning + Acting):}
    \begin{itemize}
        \item Sau khi tổng hợp và kiểm tra an toàn, hệ thống tự đánh giá chất lượng phản hồi.
        \item Nếu phát hiện thiếu bằng chứng hoặc độ tin cậy thấp, bộ điều khiển ReAct quyết định: truy xuất thêm dữ liệu, chạy lại tác tử, hoặc yêu cầu người dùng làm rõ câu hỏi.
        \item Quá trình lặp tối đa 2 vòng, đảm bảo phản hồi cuối cùng đạt chất lượng cao nhất có thể.
    \end{itemize}

    \item \textbf{Bảng đen liên tác tử (Agent Blackboard):}
    \begin{itemize}
        \item Các tác tử có thể đọc và ghi vào một không gian chia sẻ chung (blackboard), cho phép tác tử này tận dụng kết quả phân tích của tác tử khác.
        \item Ví dụ: Investment Advisor có thể đọc kết quả từ Market Analysis và Legal Advisor để đưa ra đánh giá toàn diện hơn.
        \item Mỗi mục trên bảng đen được gán nhãn tác giả, loại nội dung, danh sách bằng chứng và độ tin cậy.
    \end{itemize}

    \item \textbf{Tổng hợp và phản hồi (Synthesis):}
    \begin{itemize}
        \item Khi nhiều agent cùng được kích hoạt, hệ thống tổng hợp kết quả thành một phản hồi duy nhất, mạch lạc, loại bỏ trùng lặp.
        \item Mỗi phản hồi đều được kiểm tra an toàn (safety validation): phải có khuyến cáo pháp lý cho Legal Advisor, khuyến cáo tài chính cho Investment Advisor.
        \item Nếu không có đủ bằng chứng, hệ thống trung thực thông báo ``Chưa có đủ thông tin'' thay vì bịa đặt.
    \end{itemize}

    \item \textbf{Ghi nhớ và cá nhân hóa (Memory System):}
    \begin{itemize}
        \item Chatbot tự động trích xuất sở thích người dùng từ hội thoại (thành phố, quận, loại hình, ngân sách) và tạo MemoryProposal.
        \item Các sở thích có độ tin cậy cao được tự động áp dụng. Các sở thích khác cần người dùng xác nhận qua API.
        \item Sở thích đã xác nhận được sử dụng để cá nhân hóa phản hồi trong các phiên chat sau.
    \end{itemize}

    \item \textbf{Đánh giá chất lượng tự động (LLM Judge):}
    \begin{itemize}
        \item Hệ thống tích hợp ``giám khảo'' LLM sử dụng Gemini để chấm điểm mỗi phản hồi theo năm tiêu chí: tính căn cứ (groundedness), tính hữu ích (helpfulness), chất lượng trích dẫn (citation\_quality), an toàn (safety), và tính đầy đủ của vết xử lý (trace\_completeness).
        \item Kết quả đánh giá được lưu vào bảng EvalRun để phân tích và cải thiện hệ thống liên tục.
    \end{itemize}
\end{enumerate}

\subsubsection{Nhóm chức năng quản trị và tự động hóa}

\begin{enumerate}
    \item \textbf{Thu thập dữ liệu tự động:}
    \begin{itemize}
        \item \textbf{Crawl tin đăng hàng ngày:} Tự động thu thập tin bán và cho thuê mới từ batdongsan.com.vn vào 02:00 giờ Hồ Chí Minh. Cơ chế watermark đảm bảo chỉ crawl dữ liệu mới, tránh trùng lặp.
        \item \textbf{Crawl dự án hàng tuần:} Thu thập thông tin dự án bất động sản vào Chủ nhật hàng tuần.
        \item \textbf{Crawl tin tức hàng tuần:} Thu thập bài viết phân tích thị trường, tin tức bất động sản.
        \item \textbf{Cập nhật kiến thức pháp lý hàng tháng:} Xử lý lại toàn bộ tài liệu pháp lý, bao gồm cả HTML và PDF.
    \end{itemize}

    \item \textbf{Xử lý dữ liệu (ETL Pipeline):}
    \begin{itemize}
        \item \textbf{Làm sạch:} Chuẩn hóa giá (tỷ, triệu, triệu/tháng), diện tích (m\textsuperscript{2}), phân loại loại hình bất động sản, chuẩn hóa địa chỉ.
        \item \textbf{Làm giàu:} Geocoding -- chuyển địa chỉ thành tọa độ (latitude, longitude). Intent tagging -- tự động gán nhãn nhu cầu dựa trên nội dung mô tả.
        \item \textbf{Chunking:} Chia văn bản thành các đoạn ngữ nghĩa (400 token, 80 token chồng lấn) với 4 loại: overview, description, location, intent\_tags.
        \item \textbf{Embedding:} Tạo vector embedding 1024 chiều cho từng chunk sử dụng mô hình BGE-M3 với batch size 16.
        \item \textbf{Nạp dữ liệu:} Upsert vào PostgreSQL qua các mô-đun chuyên biệt cho từng loại dữ liệu.
    \end{itemize}

    \item \textbf{Quản trị và giám sát hệ thống:}
    \begin{itemize}
        \item \textbf{Admin Dashboard:} Giao diện quản trị với các tab: Pipeline Readiness (trạng thái các nguồn dữ liệu), Agent Traces (lịch sử xử lý với latency và status), Chat Feedback (phản hồi người dùng), Memory Proposals (các đề xuất đang chờ xác nhận).
        \item \textbf{Observability toàn diện:} Hệ thống tracing ghi lại toàn bộ luồng xử lý: intent, agents được chọn, thời gian từng bước, kết quả retrieval, LLM calls, errors. Mỗi request chatbot được lưu thành một AgentTrace với đầy đủ các AgentTraceStep.
        \item \textbf{Theo dõi chi phí LLM:} Hệ thống theo dõi số token sử dụng và chi phí ước tính cho mỗi lần gọi Gemini, có cơ chế ngân sách hàng tháng (AGENT\_LLM\_MONTHLY\_BUDGET\_USD) để kiểm soát chi phí.
        \item \textbf{Monitoring:} Prometheus + Grafana + Alertmanager cho giám sát hiệu năng và cảnh báo.
    \end{itemize}
\end{enumerate}

\subsection{Yêu cầu phi chức năng}

Bảng \ref{tab:nfr} tổng hợp các yêu cầu phi chức năng của hệ thống.

\begin{longtable}{|p{4cm}|p{9.5cm}|}
    \hline
    \textbf{Tiêu chí} & \textbf{Yêu cầu chi tiết} \\
    \hline
    Hiệu năng & 
    \begin{itemize}
        \item API listings phản hồi $<500$ms với bộ lọc cơ bản, $<1000$ms với bộ lọc phức tạp
        \item Chatbot pipeline hoàn thành $<10$s cho truy vấn thông thường, $<45$s cho truy vấn phức tạp đa tác tử
        \item Trang web tải trong $<3$s (First Contentful Paint) trên kết nối 4G
    \end{itemize} \\
    \hline
    Khả năng mở rộng &
    \begin{itemize}
        \item Kiến trúc microservices: 6 dịch vụ độc lập, mỗi service có thể scale riêng
        \item Database connection pool: 20 connections + 10 overflow
        \item Redis caching giảm tải database cho truy vấn lặp lại
        \item Airflow DAGs hỗ trợ parallel task execution
    \end{itemize} \\
    \hline
    Bảo mật &
    \begin{itemize}
        \item Xác thực JWT với HS256, token hết hạn sau 24h
        \item Mật khẩu hash bằng bcrypt với salt
        \item Chống SQL injection qua SQLAlchemy ORM (parameterized queries)
        \item Chống XSS qua Next.js auto-escaping
        \item CORS whitelist qua biến môi trường CORS\_ORIGINS
        \item Internal API giữa các service bảo vệ bằng AGENT\_INTERNAL\_KEY
    \end{itemize} \\
    \hline
    Độ tin cậy &
    \begin{itemize}
        \item Fallback mechanism 3 cấp: Agent Service $\rightarrow$ internal pipeline $\rightarrow$ thông báo an toàn
        \item Retry policy cho embedding: 3 lần với exponential backoff
        \item Retry policy cho Airflow tasks: 3 lần, 5--30 phút
        \item Health check endpoints cho tất cả services
        \item Ready check với 30s cache TTL trước khi truy xuất dữ liệu
    \end{itemize} \\
    \hline
    Khả năng bảo trì &
    \begin{itemize}
        \item Mã nguồn phân chia module rõ ràng theo domain
        \item Docstring tiếng Anh cho tất cả hàm/class
        \item Type hints đầy đủ (Python) và TypeScript (frontend)
        \item Hơn 60 file test bao phủ toàn bộ module
        \item ESLint cho frontend, compileall cho Python
    \end{itemize} \\
    \hline
    Tương thích ngôn ngữ &
    \begin{itemize}
        \item Giao diện người dùng: 100\% tiếng Việt
        \item Chatbot: hiểu và phản hồi bằng tiếng Việt
        \item Embedding: BGE-M3 tối ưu cho hơn 100 ngôn ngữ bao gồm tiếng Việt
        \item URL slugs: tiếng Việt không dấu
    \end{itemize} \\
    \hline
    \caption{Các yêu cầu phi chức năng của hệ thống}
    \label{tab:nfr}
\end{longtable}

\section{Thiết kế kiến trúc hệ thống}

\subsection{Kiến trúc tổng thể}

Hệ thống được thiết kế theo kiến trúc microservices với sáu dịch vụ độc lập, mỗi dịch vụ đảm nhận một trách nhiệm riêng biệt và giao tiếp với nhau qua REST API nội bộ. Hình~\ref{fig:architecture} thể hiện kiến trúc tổng thể.

\begin{figure}[H]
    \centering
    \includegraphics[angle=90,width=0.95\textheight]{figures/architecturepng.png}
    \caption{Kiến trúc tổng thể hệ thống Nền tảng Tư vấn Bất động sản Tích hợp Chatbot}
    \label{fig:architecture}
\end{figure}

\subsubsection{Tầng Người dùng (User Layer)}

Người dùng tương tác với hệ thống qua trình duyệt web (desktop và mobile). Có hai nhóm người dùng chính:
\begin{itemize}
    \item \textbf{Người dùng cuối:} Truy cập các trang duyệt bất động sản, xem thống kê thị trường và tương tác với chatbot tư vấn đa tác tử.
    \item \textbf{Quản trị viên:} Truy cập admin dashboard để giám sát pipeline, theo dõi agent traces và quản lý phản hồi người dùng.
\end{itemize}

\subsubsection{Tầng Frontend -- Next.js 16 (Cổng 3000)}

Frontend được xây dựng với Next.js 16 và React 19, sử dụng App Router để tổ chức các trang. Các đặc điểm thiết kế chính:
\begin{itemize}
    \item \textbf{Server-Side Rendering (SSR):} Các trang danh sách và chi tiết bất động sản được render phía server, đảm bảo SEO và thời gian tải nhanh.
    \item \textbf{React Server Components (RSC):} Mặc định sử dụng Server Components để giảm JavaScript gửi về client.
    \item \textbf{Tailwind CSS v4:} Utility-first CSS framework đảm bảo giao diện nhất quán, responsive trên mọi thiết bị.
    \item \textbf{Typed API Client:} File \texttt{lib/api.ts} cung cấp các hàm gọi backend với TypeScript type safety, tự động gắn JWT token vào header.
    \item \textbf{Chat Widget:} Thành phần nổi (floating) luôn hiển thị, cho phép người dùng tương tác với chatbot ở bất kỳ trang nào.
\end{itemize}

Frontend chỉ giao tiếp trực tiếp với Backend API qua REST. Frontend không được phép gọi trực tiếp Agent Service hay Pipeline Worker -- đây là nguyên tắc thiết kế quan trọng để đảm bảo bảo mật và khả năng bảo trì.

\subsubsection{Tầng Backend API -- FastAPI (Cổng 8000)}

Backend API là cổng giao tiếp chính của toàn bộ hệ thống, đảm nhận mọi request từ frontend. Backend được tổ chức thành 4 lớp:

\begin{enumerate}
    \item \textbf{Lớp Models (ORM):} Định nghĩa 16+ bảng dữ liệu với SQLAlchemy 2.0 async. Mỗi model ánh xạ trực tiếp đến một bảng trong PostgreSQL. Các mối quan hệ (relationships) được định nghĩa rõ ràng: Listing $\rightarrow$ Chunk (polymorphic), User $\rightarrow$ ChatSession $\rightarrow$ ChatMessage, User $\rightarrow$ UserPreference.

    \item \textbf{Lớp Schemas (Pydantic v2):} Định nghĩa contract cho request/response. Pydantic v2 tự động validate dữ liệu đầu vào, sinh JSON Schema và tài liệu OpenAPI. Các schema được thiết kế tách biệt cho từng nghiệp vụ.

    \item \textbf{Lớp Routers (API Endpoints):} Xử lý yêu cầu HTTP, gọi dịch vụ tương ứng và trả về phản hồi. Máy chủ có 9 mô-đun định tuyến: listings, market, projects, articles, chat, auth, preferences, metrics, admin. Các bộ định tuyến sử dụng dependency injection để nhận database session và xác thực người dùng.

    \item \textbf{Lớp Services (Business Logic):} Chứa toàn bộ logic nghiệp vụ, được chia thành 3 nhóm chính:
    \begin{itemize}
        \item \textbf{Chatbot Pipeline (\texttt{services/chatbot/}):} Pipeline multi-agent nội bộ với router (phân loại ý định), orchestrator (điều phối và tổng hợp), 4 agent chuyên biệt, context builder, memory manager. Đây là fallback khi Agent Service không khả dụng.
        \item \textbf{RAG Services (\texttt{services/rag/}):} Các dịch vụ truy xuất thông tin: hybrid search (kết hợp SQL + vector + rerank + cache), embeddings (BGE-M3 wrapper với batch processing và retry), cache (Redis caching cho embedding và kết quả tìm kiếm).
        \item \textbf{Agent Service Client (\texttt{services/agent\_service/}):} HTTP client giao tiếp với Agent Service qua internal API, sử dụng internal key authentication. Bao gồm observability module để ghi nhận AgentTrace, AgentTraceStep, EvalRun vào database.
    \end{itemize}
\end{enumerate}

\subsubsection{Tầng Agent Service -- Trái tim AI (Cổng 8100)}

Agent Service là dịch vụ vi mô quan trọng nhất của hệ thống, được thiết kế để xử lý các tác vụ trí tuệ nhân tạo phức tạp thông qua đồ thị trạng thái LangGraph. Dịch vụ này hoạt động độc lập với Backend API, giao tiếp qua REST với xác thực khóa nội bộ, đảm bảo isolation hoàn toàn về tài nguyên và lỗi.

\textbf{Kiến trúc đồ thị trạng thái:}

Đồ thị trạng thái LangGraph gồm 14 nút được tổ chức thành bốn pha xử lý chính, minh họa trong Hình~\ref{fig:agent-graph}.

\begin{figure}[H]
    \centering
    \begin{tabular}{|c|c|c|c|}
        \hline
        \multicolumn{4}{|c|}{\textbf{Đồ thị Trạng thái LangGraph -- 14 Nút, 4 Pha}} \\
        \hline
        \textbf{Pha 1: Hiểu} & \textbf{Pha 2: Truy xuất} & \textbf{Pha 3: Suy luận} & \textbf{Pha 4: Phản hồi} \\
        \hline
        \textbf{1.} context\_builder & \textbf{4.} query\_understanding & \textbf{6.} specialist\_agents & \textbf{9.} synthesizer \\
        \quad -- Chuẩn hóa query & \quad -- Viết lại query & \quad -- 6 agent song song & \quad -- Tổng hợp kết quả \\
        \quad -- Ngữ cảnh hội thoại & \quad -- Trích xuất bộ lọc & \textbf{7.} investment\_model & \textbf{10.} safety\_validator \\
        \hline
        \textbf{2.} readiness\_checker & \textbf{5.} retrieval\_planner & \quad -- Phân tích tài chính & \quad -- Khuyến cáo pháp lý \\
        \quad -- Kiểm tra nguồn DL & \quad -- Lập kế hoạch & \textbf{8.} committee\_review & \quad -- Khuyến cáo tài chính \\
        \quad -- Cache 30s TTL & \quad -- Hybrid search & \quad -- Đánh giá hội đồng & \quad -- Xác thực bằng chứng \\
        \hline
        \textbf{3.} router & & & \textbf{11.} react\_controller \\
        \quad -- Keyword/LLM/Hybrid & & & \quad -- Quyết định lặp \\
        \quad -- Chọn agent mục tiêu & & & \textbf{12.} react\_retrieval \\
        \textbf{3a.} clarification & & & \quad -- Truy xuất bổ sung \\
        \quad -- Làm rõ câu hỏi & & & \textbf{13.} memory\_proposals \\
        & & & \quad -- Đề xuất sở thích \\
        \hline
    \end{tabular}
    \caption{Tổng quan 14 nút trong đồ thị trạng thái LangGraph của Agent Service}
    \label{fig:agent-graph}
\end{figure}

\newpage

\textbf{Luồng xử lý chi tiết từng pha:}

\textit{Pha 1 -- Hiểu (Understand):}

\begin{enumerate}[leftmargin=*]
    \item \textbf{context\_builder} -- Nhận AgentChatRequest từ Backend, chuẩn hóa truy vấn (strip accents, lowercase), compact lịch sử hội thoại (tối đa 5 lượt gần nhất), trích xuất sở thích người dùng từ UserPreference. Kết quả được đưa vào AgentGraphState để các nút sau sử dụng.

    \item \textbf{readiness\_checker} -- Song song kiểm tra 4 nguồn dữ liệu (listings, projects, news, legal) qua SQL count. Mỗi nguồn được đánh giá ``ready'' nếu có cả parent records và chunk records. Kết quả được cache 30 giây để tránh query lặp lại. Nguồn không sẵn sàng sẽ bị loại khỏi retrieval plan.

    \item \textbf{router} -- Phân loại ý định và chọn danh sách tác tử. Hỗ trợ ba chế độ: \texttt{rule} (khớp từ khóa tiếng Việt không dấu với KEYWORDS\_BY\_AGENT), \texttt{llm} (gọi Gemini 2.0 Flash với prompt thiết kế riêng), \texttt{hybrid} (kết hợp cả hai). Router cũng phát hiện truy vấn thiếu thông tin và kích hoạt nhánh clarification.

    \item \textbf{clarification} -- Khi router phát hiện truy vấn không đủ thông tin (ví dụ: ``tìm nhà'' không có khu vực, ngân sách), nút này tạo câu hỏi làm rõ và kết thúc sớm đồ thị, trả về clarifying\_question cho người dùng.
\end{enumerate}

\textit{Pha 2 -- Truy xuất (Retrieve):}

\begin{enumerate}[leftmargin=4em]
    \item[5.] \textbf{query\_understanding} -- Phân tích sâu ngữ nghĩa truy vấn. Ở chế độ rule: trích xuất bộ lọc có cấu trúc bằng regex (listing\_type, property\_type, city, district, min/max price, min/max area, bedrooms). Ở chế độ LLM (khi AGENT\_QUERY\_REWRITE\_ENABLED): gọi Gemini để viết lại truy vấn, mở rộng truy vấn (tối đa 3 biến thể), và suy luận bộ lọc ẩn. Bộ lọc được validate qua ALLOWED\_FILTERS (18 trường được phép).

    \item[6.] \textbf{retrieval\_planner} -- Lập kế hoạch truy xuất dựa trên ba yếu tố: (1) agents được chọn cần những domain nào (AGENT\_RETRIEVAL\_DOMAINS mapping), (2) nguồn dữ liệu nào đang sẵn sàng (từ readiness\_checker), (3) bộ lọc đã trích xuất. Mỗi retrieval task được định nghĩa với: task\_id, domain, tool (search\_listings/search\_projects/search\_articles/lookup\_market\_metrics), query, filters, top\_k (mặc định 20), rerank\_top\_k (mặc định 5). Các task có dependency được giải quyết (ví dụ: investment\_advisor phụ thuộc vào market data). Task được thực thi song song qua asyncio.gather, kết quả được chuẩn hóa thành Evidence objects với source\_identity, facts, matched\_chunks.
\end{enumerate}

\textit{Pha 3 -- Suy luận (Reason):}

\begin{enumerate}[leftmargin=4em]
    \item[7.] \textbf{specialist\_agents} -- Thực thi song song các tác tử được chọn từ router. Mỗi tác tử nhận evidence được gán riêng và sinh phản hồi chuyên biệt. Ở chế độ deterministic (mặc định): mỗi agent sử dụng template Python để định dạng kết quả từ evidence -- mô tả từng bất động sản kèm giá, diện tích, vị trí và nguồn tham khảo. Ở chế độ LLM (khi AGENT\_SPECIALIST\_LLM\_ENABLED): mỗi agent gọi Gemini với system prompt riêng và evidence context, cho phản hồi tự nhiên và có chiều sâu hơn. Kết quả mỗi agent được chuẩn hóa thành SpecialistResult với status, content, evidence\_ids\_used, confidence, warnings, sources. Các agent đồng thời ghi vào Agent Blackboard -- không gian chia sẻ cho phép các agent đọc kết quả của nhau.

    \item[8.] \textbf{investment\_model} -- Chỉ chạy khi investment\_advisor được kích hoạt. Xây dựng mô hình đầu tư từ evidence đa miền: property evidence (giá, diện tích, vị trí), market evidence (giá thuê trung bình, xu hướng), legal evidence (tình trạng pháp lý), project evidence (thông tin chủ đầu tư), news evidence (tin tức liên quan). Tính toán các chỉ số tài chính với giả định mặc định (equity\_ratio=0.4, loan\_ratio=0.6, interest\_rate\_annual=0.1, loan\_term\_years=20, vacancy\_months=1, operating\_cost\_ratio=0.08): giá trị đầu tư, dòng tiền hàng năm, NPV, ROI. Kết quả được cấu trúc thành investment\_case (property\_summary, market\_summary, legal\_summary, project\_summary, news\_summary, missing\_evidence) và investment\_metrics (từng chỉ số kèm value và unit).

    \item[9.] \textbf{committee\_review} -- Đánh giá investment\_case từ ba góc nhìn độc lập dựa trên cùng bộ evidence từ agent\_blackboard:
    \begin{itemize}
        \item \textbf{Bull (Lạc quan):} Nhấn mạnh tiềm năng tăng giá, vị trí đắc địa, tiện ích xung quanh, chủ đầu tư uy tín. Mức độ rủi ro: thấp.
        \item \textbf{Bear (Thận trọng):} Nhấn mạnh rủi ro thanh khoản, chi phí vốn cao, biến động thị trường, rủi ro pháp lý. Mức độ rủi ro: cao.
        \item \textbf{Skeptic (Hoài nghi):} Đặt câu hỏi về các giả định, yêu cầu thêm bằng chứng, chỉ ra các yếu tố chưa được xác minh. Mức độ rủi ro: trung bình.
    \end{itemize}
    Mỗi góc nhìn đưa ra stance, claims (kèm evidence\_ids), confidence, risk\_level, và suggested\_actions. Kết quả hội đồng được đưa vào synthesizer để tổng hợp cùng với các agent outputs khác.
\end{enumerate}

\textit{Pha 4 -- Phản hồi (Deliver):}

\begin{enumerate}[leftmargin=4em]
    \item[10.] \textbf{synthesizer} -- Tổng hợp tất cả kết quả từ các agent, investment\_model và committee\_review thành một phản hồi cuối cùng. Quy trình: (1) thu thập tất cả agent outputs; (2) deduplicate warnings theo stable key (code+domain+message); (3) deduplicate sources theo type+id; (4) validate evidence references -- mọi evidence được trích dẫn phải tồn tại trong evidence\_for\_agent; (5) nếu không có evidence nào, fallback ``Chưa có đủ thông tin để trả lời câu hỏi này''; (6) nếu AGENT\_SPECIALIST\_LLM\_ENABLED, gọi Gemini synthesis prompt để tạo phản hồi mạch lạc; nếu không, ghép nối deterministic outputs. Kết quả bao gồm final\_response, suggested\_actions, sources, claims (mỗi claim kèm evidence\_ids).

    \item[11.] \textbf{safety\_validator} -- Kiểm tra an toàn phản hồi trước khi gửi đến người dùng. Các quy tắc: (1) nếu legal\_advisor được chạy nhưng phản hồi không chứa legal disclaimer, thêm cảnh báo ``legal\_disclaimer\_missing''; (2) nếu investment\_advisor được chạy nhưng phản hồi không chứa financial disclaimer, thêm cảnh báo ``financial\_disclaimer\_missing''; (3) các agent có nguồn dữ liệu (SOURCE\_BACKED\_AGENTS: legal\_advisor, news\_agent, project\_agent, property\_search) phải có ít nhất một evidence tham chiếu, nếu không thêm cảnh báo ``agent\_answer\_missing\_valid\_evidence''. Cảnh báo được thêm vào state.warnings và hiển thị trong trace.

    \item[12.] \textbf{react\_controller} -- Điều khiển vòng lặp phản ứng. Chỉ hoạt động khi AGENT\_REACT\_ENABLED=true. Kiểm tra state.warnings: nếu có cảnh báo nghiêm trọng (final\_response\_missing\_sources, agent\_answer\_missing\_valid\_evidence), quyết định hành động tiếp theo: \texttt{retrieve\_more} (truy xuất thêm với bộ lọc mở rộng), \texttt{run\_specialist} (chạy lại specialist agents), \texttt{ask\_clarification} (yêu cầu người dùng làm rõ), hoặc \texttt{finalize} (chấp nhận kết quả hiện tại). Số vòng lặp tối đa: AGENT\_REACT\_MAX\_ITERATIONS (mặc định 2).

    \item[13.] \textbf{react\_retrieval} -- Thực hiện truy xuất bổ sung khi react\_controller yêu cầu. Sử dụng filters từ react\_decision để mở rộng hoặc điều chỉnh retrieval plan. Kết quả được merge vào evidence\_by\_id và evidence\_for\_agent hiện có, sau đó quay lại specialist\_agents.

    \item[14.] \textbf{memory\_proposals} -- Trích xuất sở thích người dùng từ query và filters đã xác định. Mapping FILTER\_TO\_MEMORY\_KEY: city $\rightarrow$ preferred\_city, district $\rightarrow$ preferred\_district, property\_type $\rightarrow$ preferred\_property\_type, listing\_type $\rightarrow$ listing\_type, bedrooms $\rightarrow$ bedrooms, max\_price $\rightarrow$ max\_budget, min\_price $\rightarrow$ min\_budget. Mỗi proposal có confidence (0.72--0.82) và requires\_user\_confirmation=True. Proposal được trả về trong AgentChatResponse để Backend lưu vào bảng memory\_proposals và hiển thị cho người dùng xác nhận.
\end{enumerate}

\textbf{Cấu trúc dữ liệu trạng thái (AgentGraphState):}

Toàn bộ trạng thái của đồ thị được lưu trong AgentGraphState -- một TypedDict với hơn 40 trường, bao gồm:

\begin{itemize}
    \item \textbf{Đầu vào:} request (AgentChatRequest với message, session\_id, user\_id, conversation\_context, preferences).
    \item \textbf{Trung gian:} normalized\_query, intent, agents\_to\_run, routing\_filters, query\_understanding, readiness, retrieval\_plan, evidence\_by\_id, evidence\_for\_agent, agent\_results, agent\_blackboard, investment\_case, investment\_assumptions, investment\_metrics, committee\_review.
    \item \textbf{Điều khiển:} needs\_clarification, force\_deterministic, react\_iteration, react\_decision, react\_actions.
    \item \textbf{Đầu ra:} final\_response, sources (list[AgentSource]), suggested\_actions, memory\_proposals, trace\_steps, warnings.
\end{itemize}

\textbf{Cơ chế đặc biệt:}

\textit{Agent Blackboard:} Không gian chia sẻ giữa các tác tử, triển khai qua blackboard dictionary trong AgentGraphState. Mỗi entry có: id, author (tên agent), type (specialist\_result, market\_insight, risk\_assessment,\dots), content, evidence\_ids, confidence (low/medium/high), created\_at\_step. Các agent có thể đọc entry của nhau qua entries\_by\_author(). Investment Advisor đặc biệt tận dụng blackboard để tổng hợp insight từ Market Analysis và Legal Advisor.

\textit{Evidence Normalization:} Mọi kết quả retrieval được chuẩn hóa thành Evidence objects với: evidence\_id (duy nhất), retrieval\_task\_id (nguồn gốc), domain (property/project/news/legal/market), source\_type, source\_identity (stable key type:id), record (dữ liệu gốc), facts (title, price, location,\dots), source (AgentSource với type, id, title, url, score), matched\_chunks (kèm vector\_distance, rerank\_score, final\_score), retrieved\_for (danh sách agent), assigned\_to (danh sách agent được phân bổ), warnings.

\textit{Cost Tracking:} Mỗi lần gọi Gemini được ghi nhận: input\_tokens, output\_tokens, estimated\_cost\_usd. Hệ thống kiểm tra AGENT\_LLM\_MONTHLY\_BUDGET\_USD trước mỗi lần gọi. Nếu vượt ngân sách, LLM bị vô hiệu hóa và hệ thống tự động chuyển sang deterministic mode.

\textit{LLM Judge:} Endpoint \texttt{/internal/agent/evaluate} cho phép đánh giá chất lượng phản hồi qua 5 metrics: groundedness (phản hồi có được hỗ trợ bởi nguồn không), helpfulness (có trực tiếp giúp người dùng không), citation\_quality (nguồn có liên quan và rõ ràng không), safety (có khẳng định pháp lý/tài chính tuyệt đối không), trace\_completeness (trace có đầy đủ agents, steps, warnings không). Judge prompt được thiết kế riêng cho tiếng Việt, yêu cầu JSON output với scores và rationales.

\subsubsection{Tầng Pipeline Worker -- Xử lý dữ liệu tự động (Cổng 8200)}

Pipeline Worker là dịch vụ vi mô chuyên trách các tác vụ xử lý dữ liệu nặng, tách biệt khỏi Backend và Agent Service để tránh ảnh hưởng đến hiệu năng phục vụ người dùng.

\textbf{Kiến trúc:} Pipeline Worker chạy FastAPI trên cổng 8200, cung cấp internal API endpoints:

\begin{itemize}
    \item \texttt{POST /internal/pipeline/crawler} -- Thực thi module crawler (sale, rent, projects, news) qua subprocess. Nhận module name, args, timeout (mặc định 7200s).
    \item \texttt{POST /internal/pipeline/ingest/listings} -- Nạp dữ liệu từ CSV vào PostgreSQL qua data\_pipeline ingestors.
    \item \texttt{POST /internal/pipeline/maintenance/cleanup-chunks} -- Dọn dẹp chunks cũ, giải phóng dung lượng.
    \item \texttt{GET /internal/pipeline/health} -- Kiểm tra trạng thái dịch vụ.
\end{itemize}

Tất cả endpoints yêu cầu xác thực X-Internal-Agent-Key. Worker chạy các tác vụ dưới dạng subprocess Python, capture stdout/stderr, và trả về kết quả đồng bộ. Các module được gọi qua \texttt{python -m <module>} với PYTHONPATH được thiết lập đúng.

\subsubsection{Tầng Dữ liệu (Data Layer)}

\textbf{PostgreSQL 16 + pgvector (Cổng 5432):}

PostgreSQL được chọn làm cơ sở dữ liệu chính nhờ khả năng lưu trữ đồng thời dữ liệu quan hệ và dữ liệu véc-tơ thông qua extension pgvector. Điều này giúp thực hiện hybrid search trong cùng một truy vấn SQL: JOIN giữa bảng listings và bảng chunks, WHERE cho bộ lọc có cấu trúc, ORDER BY cosine distance cho tìm kiếm ngữ nghĩa.

Dữ liệu quan hệ gồm 16+ bảng: listings (30+ trường), projects, articles, users, chat\_sessions, chat\_messages, chunks, agent\_traces, agent\_trace\_steps, agent\_llm\_calls, agent\_retrieval\_events, eval\_runs, pipeline\_runs, source\_readiness, user\_preferences, memory\_proposals, chat\_feedback.

Dữ liệu véc-tơ được lưu trong bảng chunks với cột embedding VECTOR(1024). Chỉ mục HNSW với tham số $m=16$ và $ef\_construction=64$ đảm bảo tốc độ tìm kiếm k-NN nhanh. Thiết kế đa hình (polymorphic) qua cặp cột (parent\_type, parent\_id) cho phép một bảng chunks phục vụ nhiều loại dữ liệu nguồn.

\textbf{Redis 7 (Cổng 6379):}

Redis được sử dụng làm bộ nhớ đệm ở hai cấp độ: embedding cache (lưu vector của câu truy vấn, tránh encode lại) và search result cache (lưu kết quả hybrid search cho truy vấn lặp lại).

\subsubsection{Tầng Thu thập \& Xử lý dữ liệu}

\textbf{Crawlers (Playwright + Stealth):}

Hệ thống crawler được tổ chức thành 5 module:
\begin{itemize}
    \item \texttt{crawler/core/}: Các thành phần dùng chung -- CSV writer (với dedup), base parser.
    \item \texttt{crawler/sale/}: Crawl tin bán (2 bước: crawl URLs $\rightarrow$ crawl details).
    \item \texttt{crawler/rent/}: Crawl tin cho thuê.
    \item \texttt{crawler/projects/}: Crawl thông tin dự án bất động sản.
    \item \texttt{crawler/news/}: Crawl tin tức và bài viết phân tích thị trường.
\end{itemize}

Crawler sử dụng Playwright với chế độ stealth để tránh bị phát hiện, hỗ trợ 4 workers song song. Dữ liệu được ghi ra file CSV UTF-8 BOM và deduplicate dựa trên product\_id.

\textbf{Data Pipeline (ETL + Embedding):}

Pipeline xử lý dữ liệu gồm 5 bước:
\begin{enumerate}
    \item \textbf{Clean (\texttt{clean.py}):} Chuẩn hóa giá, diện tích, phân loại loại hình, chuẩn hóa địa chỉ.
    \item \textbf{Enrich (\texttt{enrich.py}):} Geocoding (địa chỉ $\rightarrow$ tọa độ), Intent Tagging (gán nhãn nhu cầu).
    \item \textbf{Chunk (\texttt{chunk.py}):} Chia văn bản thành các đoạn ngữ nghĩa 400 token, overlap 80 token, 4 loại: overview, description, location, intent\_tags.
    \item \textbf{Embed (\texttt{embed.py}):} Tạo vector embedding 1024 chiều bằng BGE-M3, batch size 16, retry 3 lần.
    \item \textbf{Load (\texttt{load\_db.py} + ingestors/):} Upsert vào PostgreSQL qua các ingestor module chuyên biệt.
\end{enumerate}

\textbf{Apache Airflow Scheduler:}

Toàn bộ pipeline được lập lịch và giám sát bởi Apache Airflow với 4 DAGs:
\begin{itemize}
    \item \textbf{daily\_listings\_dag (02:00 HCM):} Crawl + xử lý tin đăng mới hàng ngày.
    \item \textbf{weekly\_projects\_dag (Chủ nhật 03:00):} Crawl + xử lý dự án.
    \item \textbf{weekly\_news\_dag (Chủ nhật 04:00):} Crawl + xử lý tin tức.
    \item \textbf{monthly\_legal\_kb\_dag (Ngày 1, 05:00):} Xử lý lại kho kiến thức pháp lý.
\end{itemize}

\section{Thiết kế cơ sở dữ liệu}

\subsection{Mô hình quan hệ}

Hệ thống sử dụng PostgreSQL 16 với extension pgvector. Mô hình dữ liệu được thiết kế xoay quanh 4 nhóm bảng chính:

\subsubsection{Nhóm bảng nghiệp vụ (Business Tables)}

\begin{itemize}
    \item \textbf{listings -- Bảng bất động sản:} Bảng quan trọng nhất với hơn 30 trường, lưu toàn bộ thông tin của một tin đăng. Khóa chính id (INTEGER). Mỗi listing có product\_id duy nhất từ nguồn crawl để deduplicate. Các chỉ mục trên: listing\_type, city, district, property\_type, price, area, bedrooms, post\_date, is\_active.

    \item \textbf{projects -- Bảng dự án:} Lưu thông tin dự án: tên, chủ đầu tư, vị trí, quy mô, khoảng giá, trạng thái, loại hình, tiện ích, mô tả.

    \item \textbf{articles -- Bảng bài viết:} Lưu tin tức và kiến thức pháp lý: tiêu đề, nội dung, category (news/legal), nguồn, ngày đăng, URL gốc.

    \item \textbf{users -- Bảng người dùng:} Lưu thông tin tài khoản: email (unique), mật khẩu đã hash, họ tên, số điện thoại, avatar, quyền.
\end{itemize}

\subsubsection{Nhóm bảng hội thoại (Chat Tables)}

\begin{itemize}
    \item \textbf{chat\_sessions -- Phiên hội thoại:} id UUID, liên kết đến user (có thể NULL cho anonymous), tiêu đề tự động sinh.
    \item \textbf{chat\_messages -- Tin nhắn:} Mỗi tin nhắn thuộc về một session, role (user/assistant/system), nội dung, agent đã xử lý, metadata JSONB (sources, citations, search context).
\end{itemize}

\subsubsection{Nhóm bảng vector và tìm kiếm (Vector \& Retrieval)}

\begin{itemize}
    \item \textbf{chunks -- Đoạn văn bản đã embedding:} Thiết kế polymorphic với parent\_type (listing/project/article) và parent\_id. Mỗi chunk có chunk\_type (overview/description/location/intent\_tags/legal), text, embedding VECTOR(1024). HNSW index trên cột embedding với cosine distance operator.
\end{itemize}

\subsubsection{Nhóm bảng giám sát \& vận hành (Observability Tables)}

\begin{itemize}
    \item \textbf{agent\_traces:} Lưu vết xử lý của mỗi request chatbot: request\_id, session\_id, intent, agents\_used, trace\_summary, full\_trace (JSONB), latency\_ms, status (success/partial/error), graph\_version, prompt\_version, model\_name.
    \item \textbf{agent\_trace\_steps:} Chi tiết từng bước trong trace: step\_name, status, latency\_ms, input/output JSONB, error.
    \item \textbf{agent\_llm\_calls:} Log mỗi lần gọi LLM: node\_name, model\_name, latency\_ms, input\_tokens, output\_tokens, estimated\_cost\_usd.
    \item \textbf{agent\_retrieval\_events:} Log mỗi sự kiện truy xuất: tool\_name, domain, filters, result\_count, latency\_ms.
    \item \textbf{eval\_runs:} Lưu kết quả đánh giá tự động: scores (5 metrics), rationales, judge\_model, graph\_version.
    \item \textbf{pipeline\_runs:} Log mỗi lần chạy Airflow DAG.
    \item \textbf{source\_readiness:} Trạng thái sẵn sàng của từng nguồn dữ liệu.
    \item \textbf{user\_preferences:} Sở thích người dùng dạng key-value JSONB.
    \item \textbf{memory\_proposals:} Đề xuất sở thích từ agent (trạng thái: pending/accepted/rejected).
    \item \textbf{chat\_feedback:} Phản hồi của người dùng về chất lượng chat (rating, comment).
\end{itemize}

\subsection{Chiến lược indexing}

\begin{itemize}
    \item \textbf{B-tree indexes cho listings:} listing\_type, city, district, property\_type, price, area, bedrooms, post\_date, is\_active -- tối ưu các truy vấn lọc và sắp xếp phổ biến.
    \item \textbf{HNSW index cho chunks:} Trên cột embedding với vector\_cosine\_ops, $m=16$, $ef\_construction=64$ -- cho phép tìm kiếm k-NN trong không gian 1024 chiều với độ phức tạp $O(\log N)$.
    \item \textbf{Composite index cho polymorphic queries:} (parent\_type, parent\_id) trên bảng chunks -- truy vấn nhanh tất cả chunks thuộc về một thực thể.
    \item \textbf{Index cho observability:} agent\_traces (session\_id, created\_at), agent\_trace\_steps (trace\_id) -- hỗ trợ truy vấn admin dashboard.
\end{itemize}

\section{Thiết kế giao diện người dùng}

Hệ thống có giao diện người dùng web responsive, tối ưu cho cả desktop và mobile, được xây dựng với Next.js 16 và Tailwind CSS v4. Các trang chính bao gồm:

\begin{itemize}
    \item \textbf{Trang chủ (\texttt{/}):} Giới thiệu nền tảng, thanh tìm kiếm nổi bật, các danh mục phổ biến, và chỉ số thị trường tổng quan.
    \item \textbf{Nhà đất bán (\texttt{/nha-dat-ban}):} Danh sách bất động sản bán dạng lưới thẻ với bộ lọc nâng cao và phân trang.
    \item \textbf{Nhà đất cho thuê (\texttt{/nha-dat-cho-thue}):} Tương tự nhưng cho tin cho thuê.
    \item \textbf{Thị trường (\texttt{/thi-truong}):} Biểu đồ tương tác hiển thị thống kê giá, phân bố loại hình, top khu vực.
    \item \textbf{Dự án (\texttt{/du-an}):} Danh sách dự án bất động sản với bộ lọc.
    \item \textbf{Tin tức (\texttt{/tin-tuc}):} Danh sách bài viết tin tức và phân tích thị trường.
    \item \textbf{Đăng nhập (\texttt{/dang-nhap}) / Đăng ký (\texttt{/dang-ky}):} Form xác thực người dùng.
    \item \textbf{Admin Dashboard (\texttt{/admin}):} Giao diện quản trị với các tab: Pipeline Readiness, Agent Traces, Chat Feedback, Memory Proposals.
\end{itemize}

Thành phần \textbf{Chat Widget} là một floating button luôn hiển thị ở góc phải dưới mọi trang. Khi click, widget mở rộng thành cửa sổ chat với: ô nhập liệu, lịch sử hội thoại (scrollable), hiển thị sources dưới dạng accordion, nút đánh giá (thích/không thích), và suggested actions. Widget gọi API \texttt{/api/v1/chat} qua typed client \texttt{lib/api.ts} với JWT token tự động đính kèm.

\newpage
